\section{Coherent finite presentations}
\label{sec:refinement}

The existing packets use finite observations and explicit error contracts.
For algebraic constructions there is a useful further organizing principle:
successive finite observations can be \emph{exact restrictions} of one
another. This avoids treating all infinite objects as analytic Cauchy
sequences and exposes the algebra actually needed by the proof.

\subsection{A presentation at each requested precision}

For fixed numerical parameters \(x\), let \(A_n\) be a presentation of
finite data at precision \(n\). In the arithmetic theory this means a
bounded validity test \(V(x,n,z)\), a bounded equivalence test
\(z\simeq_{x,n}z'\), and PR restriction operations
\(\pi_{m,n}\) for \(n\le m\). Their laws, on valid inputs, are
\[
 \pi_{n,n}(z)\simeq_{x,n}z,\qquad
 \pi_{k,n}(\pi_{m,k}(z))\simeq_{x,n}\pi_{m,n}(z)
 \quad(n\le k\le m).
\]
All operations respect the equivalence tests and preserve validity.
For polynomials modulo \(T^{n+1}\), restriction deletes high
coefficients. For residues modulo \(q^n\), it reduces to the smaller
modulus. For finite matrix identities, precision may instead mean the
finite list of entries or identities requested.

A represented object is a specified PR producer \(a(x,n)\), together
with proofs of its validity and coherence
\[
 \pi_{m,n}(a(x,m))\simeq_{x,n}a(x,n)\qquad(n\le m).
\]
The first-order theorem quantifies over the numerical parameters and the
requested precisions. It does not assert the existence of a function
object inside \ipra. One can interpret the family externally as an
element of an inverse limit, but that interpretation is explanatory, not
an additional axiom used to establish the finite claims.

\begin{theorem}[Finite uniqueness supplies coherence]
\label{thm:coherence}
Let \(C(x,n,z)\) be a bounded specification that includes validity at
precision \(n\). Suppose a fixed PR producer
\(a(x,n)\) satisfies it, restrictions preserve it, and it determines a
unique equivalence class at every precision:
\begin{align*}
 &C(x,n,a(x,n)),\\
 &n\le m\land C(x,m,z)\to C(x,n,\pi_{m,n}(z)),\\
 &C(x,n,z)\land C(x,n,z')\to z\simeq_{x,n}z'.
\end{align*}
If these implications have \ipra\ derivations, so does the coherence
formula for \(a\).
\end{theorem}
\begin{proof}
Fix \(x,n,m\) with \(n\le m\). Correctness at \(m\) and restriction
give \(C(x,n,\pi_{m,n}(a(x,m)))\). Correctness at \(n\) gives
\(C(x,n,a(x,n))\). Apply finite uniqueness. No induction and no choice
principle is required in this final step.
\end{proof}

The statement deliberately includes existence \emph{with a fixed PR
producer}. A bare theorem \(\forall n\,\exists z\,C(n,z)\), even with
uniqueness, does not by itself provide a PR function. The lifting theorem
in the next chapter constructs the producer before coherence is invoked.
Moreover the specification includes the selected initial residue or
constant term. Without that seed there may be several roots and the
uniqueness hypothesis would be false.

\subsection{Operations carry an explicit precision demand}

An operation producing an output observation at level \(n\) must state
which input levels it reads. For product of ordinary jets the demand is
\(n\) for each input. For a derivative it is \(n+1\). For a coefficient
in an implicit solution the demand is the polynomial data and the seed
plus its finite inverse certificate. Once these demands are stated,
composition is ordinary composition of their PR bounds and programs.

Suppose a first operation has witness \(f(x)\) and proof
\(P(x)\to Q(x,f(x))\), and a second operation has witness \(g(x,y)\)
and proof \(Q(x,y)\to R(x,y,g(x,y))\). Substitution and implication
elimination give
\[
 P(x)\to R(x,f(x),g(x,f(x))).
\]
This is a syntactic proof transformation, not an appeal to semantic
soundness. The module \texttt{ObjectComposition.lean} implements it for
the repository's actual \texttt{Prf PRAAx} relation, including capture-free
substitution and specialization of parameters.

Coherent operations can similarly be composed. Restrict the outer result,
use its restriction law, apply the restriction laws of its arguments,
and then use congruence. Every equality is a finite equality at a stated
level. No equality test for the entire infinite object is needed.

\subsection{Three kinds of precision, with different meanings}

\begin{center}\small
\begin{tabular}{P{.20\linewidth}P{.30\linewidth}P{.40\linewidth}}
\toprule
Presentation & Finite observation & What passing every level means externally\\
\midrule
Ordinary jets & Coefficients through degree \(n\), exact modulo
\(T^{n+1}\) & A formal power series identity; it does not supply a radius
of analytic convergence.\\[3pt]
Residue towers & A root modulo \(q^n\), with a selected residue modulo
\(q\) & A compatible \(q\)-adic solution; it need not be a rational or
integer solution.\\[3pt]
Metric names & Rational observations with explicit error bounds & A real
or complex object in a stated represented family; exact zero tests do not
follow from arbitrary finite approximations.\\
\bottomrule
\end{tabular}
\end{center}

These are different embeddings of mathematics. Their finite arithmetic
theorems are useful in their own right. A conversion between them needs a
separate proof: convergence estimates convert suitable jets to analytic
names, and a global arithmetic argument may sometimes turn local residue
information into a rational solution. Neither conversion is automatic.

\subsection{Certificates: validity, production, and scope}

For every certificate language record three separate statements:
\[
 \CertCheck(x,c)=1\to \Spec(x,c),\qquad
 \CertCheck(x,\operatorname{produce}(x))=1,
\]
and the exact input class on which the second statement is proved.
The specification itself must be an explicit arithmetic formula, not an
unexplained proof field whose contents already are the desired theorem.

In the lifting application, the initial certificate consists of a root
modulo the first filtration level and a matrix inverse there. Both are
finite checks. The algorithm then produces all requested refinements.
For rational jets a nonzero rational determinant lets finite elimination
construct the inverse. For residues, the inverse may be supplied or
found by a bounded search modulo \(q\), provided invertibility is
established. In the norm application the certificate is constructed
directly: form the multiplication matrix and its determinant. Its exact
degree is proved from the weighted leading terms. No external oracle
supplies the norm.

The later covering corollary has a different status. Given a list of
covering subspaces, it proves a point bound. Producing such a list for
the integral or rational points on an arbitrary curve remains a separate
mathematical obligation. The type of the output is kept visible precisely
so that a counting theorem cannot be mistaken for a height bound.
